% !TEX root = ../BeamerTemplate.tex
%%%%%%%%%%%%%%%%%%%%%%%%%%%%%%%%%%%%%%%%%%%%%%%%%%%%%%%%%%%%%%%%%%%%%%%%%%%%%%%%%%

\section{Monte Carlo 9x9 Position Simulation}

\subsection{Simulation Design}

\begin{frame}[t,shrink=10]{Monte Carlo Simulation Objective}
    \begin{columns}[T,onlytextwidth]
        \begin{column}{0.52\textwidth}
            \begin{block}{Question}
                How does the RX antenna lens affect a 38-GHz indoor channel
                when nine transmitters are placed at random positions?
            \end{block}
            \begin{itemize}
                \item Room size: $20\times20\times3$~m
                \item Nine randomly positioned TXs per drop
                \item Nine RX angular configurations
                \item With-lens and without-lens cases
                \item 20 Monte Carlo drops
            \end{itemize}
        \end{column}
        \begin{column}{0.44\textwidth}
            \begin{exampleblock}{Paired experiment}
                Each lens/no-lens comparison uses identical TX positions.
                This isolates the effect of the RX pattern from random
                geometry changes.
            \end{exampleblock}
            \begin{alertblock}{Evidence status}
                With only 20 drops and 100,000 path samples per source, the
                results are preliminary and should be shared with caveats.
            \end{alertblock}
        \end{column}
    \end{columns}
\end{frame}

\begin{frame}[t,shrink=15]{Lens and No-Lens Configurations}
    \begin{columns}[T,onlytextwidth]
        \begin{column}{0.47\textwidth}
            \begin{block}{Without lens}
                \begin{itemize}
                    \item Same patch RX pattern for every angle label
                    \item RX labels retained only for paired matching
                    \item Expected to provide relatively uniform angular
                    performance
                \end{itemize}
            \end{block}
        \end{column}
        \begin{column}{0.49\textwidth}
            \begin{block}{With lens}
                Different HFSS far-field patterns are used at:
                \[
                    +60^\circ,+45^\circ,+30^\circ,+15^\circ,0^\circ,
                \]
                \[
                    -15^\circ,-30^\circ,-45^\circ,-60^\circ.
                \]
                Performance can therefore depend strongly on angular
                alignment.
            \end{block}
        \end{column}
    \end{columns}
    \begin{exampleblock}{Dataset size}
        Each notebook contains $20$ drops $\times$ $9$ RX configurations
        $=180$ drop--RX results; the paired comparison contains 180 pairs.
    \end{exampleblock}
\end{frame}

\begin{frame}[t,shrink=12]{Simulation Scope and Parameters}
    \centering
    \small
    \begin{tabular}{@{}ll@{\qquad}ll@{}}
        \toprule
        \textbf{Parameter} & \textbf{Value} &
        \textbf{Parameter} & \textbf{Value} \\
        \midrule
        Carrier frequency & 38~GHz &
        Bandwidth & 400~MHz \\
        Frequency points & 401 &
        Drops & 20 \\
        TXs per drop & 9 &
        RX configurations & 9 \\
        Random seed & 20260718 &
        Path samples/source & 100,000 \\
        TX height & 1.0--2.5~m &
        Minimum TX--RX range & 1.0~m \\
        Wall margin & 0.75~m &
        Power per TX & 10~dBm \\
        Noise figure & 7~dB &
        Noise temperature & 290~K \\
        Outage threshold & 1~bit/s/Hz &
        Results/notebook & 180 \\
        \bottomrule
    \end{tabular}

    \vspace{5mm}
    \begin{alertblock}{TX pattern clarification}
        Although the notebook filenames contain \texttt{Patch}, the Monte
        Carlo section explicitly forces the TX pattern to isotropic
        (\texttt{iso}).
    \end{alertblock}
\end{frame}

\begin{frame}[t,shrink=8]{Random-Position Spatial Coverage}
    \begin{columns}[T,onlytextwidth]
        \begin{column}{0.47\textwidth}
            \begin{block}{Stored TX-coordinate coverage}
                \centering
                \begin{tabular}{@{}lc@{}}
                    \toprule
                    \textbf{Dimension} & \textbf{Range} \\
                    \midrule
                    $x$ & $-9.24$ to $9.19$~m \\
                    $y$ & $-9.24$ to $9.19$~m \\
                    $z$ & $1.01$ to $2.50$~m \\
                    \bottomrule
                \end{tabular}
            \end{block}
        \end{column}
        \begin{column}{0.49\textwidth}
            \begin{block}{TX--RX distance}
                \centering
                \begin{tabular}{@{}lc@{}}
                    \toprule
                    \textbf{Statistic} & \textbf{Distance} \\
                    \midrule
                    Minimum & 1.22~m \\
                    Mean & 10.56~m \\
                    Maximum & 19.57~m \\
                    \bottomrule
                \end{tabular}
            \end{block}
        \end{column}
    \end{columns}
    \begin{exampleblock}{Coverage check}
        The 180 stored TX points span a broad portion of the permitted room
        area and are consistent with the configured sampling constraints.
    \end{exampleblock}
\end{frame}

\subsection{Calculation Methodology}

\begin{frame}[t,shrink=18]{Channel and Capacity Calculation}
    \begin{columns}[T,onlytextwidth]
        \begin{column}{0.52\textwidth}
            \begin{block}{Non-coherent power combination}
                \[
                    P_{\mathrm{RX}}(f)=P_{\mathrm{TX}}
                    \sum_{i=1}^{9}|H_i(f)|^2
                \]
                Powers from the nine TX channels are added, not their complex
                voltages.
            \end{block}
            \begin{block}{Equivalent SISO spectral efficiency}
                \[
                    C=\mathbb E_f\!
                    \left[\log_2\!\left(1+\mathrm{SNR}(f)\right)\right]
                \]
            \end{block}
        \end{column}
        \begin{column}{0.44\textwidth}
            \begin{itemize}
                \item Noise is calculated from $kTB$ and the configured noise
                figure.
                \item Ideal throughput is $B\,C$.
                \item RMS delay spread is extracted from the IFFT-derived
                power-delay profile.
                \item Shannon capacity excludes modulation, coding,
                overhead, interference, estimation, and RF losses.
            \end{itemize}
        \end{column}
    \end{columns}
\end{frame}

\begin{frame}[t,shrink=12]{Virtual 9x9 MIMO Evaluation}
    \begin{columns}[T,onlytextwidth]
        \begin{column}{0.53\textwidth}
            \begin{block}{Construction}
                The nine RX channels associated with different lens angles
                are combined into a virtual $9\times9$ matrix.
            \end{block}
            \textbf{Reported metrics:}
            \begin{itemize}
                \item Condition number
                \item Effective rank
                \item Capacity at 10-dB SNR
            \end{itemize}
        \end{column}
        \begin{column}{0.43\textwidth}
            \begin{alertblock}{Interpretation limit}
                The nine RX states are not measured simultaneously at one
                physical receiver.
            \end{alertblock}
            \begin{exampleblock}{Correct interpretation}
                These metrics indicate virtual angular-channel diversity and
                conditioning. They are not directly realizable capacity for a
                physical single-antenna RX.
            \end{exampleblock}
        \end{column}
    \end{columns}
\end{frame}

\subsection{Aggregate Results}

\begin{frame}[t,shrink=18]{Aggregate Lens and No-Lens Results}
    \centering
    \scriptsize
    \begin{tabular}{lrrr}
        \toprule
        \textbf{Metric} & \textbf{No lens} & \textbf{Lens} &
        \textbf{Change} \\
        \midrule
        Median capacity (bit/s/Hz) & 7.29 & 7.67 &
        $+0.38$ ($+5.2\%$) \\
        5th-percentile capacity & 6.48 & 5.02 & $-1.46$ \\
        95th-percentile capacity & 9.21 & 9.68 & $+0.47$ \\
        Median throughput (Gbit/s) & 2.92 & 3.07 &
        $+0.15$ ($+5.1\%$) \\
        Median RMS delay spread (ns) & 68.29 & 65.01 &
        $-3.28$ ($-4.8\%$) \\
        Outage below 1 bit/s/Hz & 0.0\% & 0.0\% & No change \\
        \bottomrule
    \end{tabular}

    \vspace{4mm}
    \begin{columns}[T,onlytextwidth]
        \begin{column}{0.47\textwidth}
            \begin{exampleblock}{Typical performance}
                The lens improves median capacity and throughput by
                approximately \textbf{5\%}.
            \end{exampleblock}
        \end{column}
        \begin{column}{0.49\textwidth}
            \begin{alertblock}{Reliability tradeoff}
                The 5th percentile degrades, indicating greater sensitivity
                to angle and configuration.
            \end{alertblock}
        \end{column}
    \end{columns}
\end{frame}

\begin{frame}[t,shrink=12]{Median Capacity by RX Angle}
    \begin{columns}[T,onlytextwidth]
        \begin{column}{0.55\textwidth}
            \centering
            \scriptsize
            \begin{tabular}{rrrr}
                \toprule
                \textbf{Angle} & \textbf{No lens} & \textbf{Lens} &
                \textbf{Difference} \\
                \midrule
                $+60^\circ$ & 7.31 & 6.70 & $-0.61$ \\
                $+45^\circ$ & 7.31 & 7.49 & $+0.18$ \\
                $+30^\circ$ & 7.31 & 7.29 & $-0.01$ \\
                $+15^\circ$ & 7.30 & 8.02 & $+0.72$ \\
                $0^\circ$   & 7.30 & 7.68 & $+0.39$ \\
                $-15^\circ$ & 7.29 & 8.28 & $+0.99$ \\
                $-30^\circ$ & 7.29 & \textbf{8.32} & $\mathbf{+1.03}$ \\
                $-45^\circ$ & 7.29 & 8.13 & $+0.84$ \\
                $-60^\circ$ & 7.29 & 7.43 & $+0.14$ \\
                \bottomrule
            \end{tabular}
        \end{column}
        \begin{column}{0.41\textwidth}
            \begin{exampleblock}{Best region}
                The strongest median performance occurs from
                $-15^\circ$ to $-45^\circ$, led by $-30^\circ$.
            \end{exampleblock}
            \begin{alertblock}{Observed asymmetry}
                $+60^\circ$ performs worse than the no-lens case. Positive
                and negative angles cannot be assumed symmetric.
            \end{alertblock}
            \small
            Verify RX orientation, HFSS mapping, parameter filtering, and
            angle signs before treating the asymmetry as physical.
        \end{column}
    \end{columns}
\end{frame}

\begin{frame}[t,shrink=10]{Paired Checks Using Identical TX Positions}
    \begin{columns}[T,onlytextwidth]
        \begin{column}{0.46\textwidth}
            \begin{block}{Across 180 pairs}
                \centering
                \resizebox{\linewidth}{!}{%
                \begin{tabular}{@{}lr@{}}
                    \toprule
                    \textbf{Paired statistic} & \textbf{Result} \\
                    \midrule
                    Median capacity difference & $+0.40$~bit/s/Hz \\
                    Mean capacity difference & $+0.15$~bit/s/Hz \\
                    Lens wins & 106/180 ($58.9\%$) \\
                    Median SNR difference & $+1.3$~dB \\
                    Median delay-spread difference & $-4.70$~ns \\
                    \bottomrule
                \end{tabular}
                }
            \end{block}
        \end{column}
        \begin{column}{0.50\textwidth}
            \begin{block}{Lens wins by angle}
                \begin{itemize}
                    \item 15/20 drops at $+15^\circ$ and $-15^\circ$
                    \item 14/20 drops at $0^\circ$ and $-30^\circ$
                    \item 13/20 drops at $-45^\circ$
                    \item Only 8/20 drops at $+60^\circ$
                \end{itemize}
            \end{block}
            \begin{alertblock}{Conclusion}
                Lens benefit depends on angular alignment and is not a
                universal improvement.
            \end{alertblock}
        \end{column}
    \end{columns}
\end{frame}

\subsection{Virtual MIMO Results}

\begin{frame}[t,shrink=12]{Virtual MIMO Channel Conditioning}
    \centering
    \small
    \begin{tabular}{lrr}
        \toprule
        \textbf{Metric} & \textbf{No lens} & \textbf{Lens} \\
        \midrule
        Median condition number & 49.4 (33.8~dB) & 54.0 (34.6~dB) \\
        5th--95th percentile range & 27.8--44.4~dB & 29.4--40.6~dB \\
        Median effective rank & 4.09 of 9 & 4.07 of 9 \\
        Median capacity at 10~dB & 21.72~bit/s/Hz & 21.11~bit/s/Hz \\
        \bottomrule
    \end{tabular}

    \vspace{5mm}
    \begin{alertblock}{No virtual-MIMO improvement}
        The lens slightly increases the condition number, leaves effective
        rank almost unchanged, and reduces median 10-dB virtual-MIMO capacity
        by approximately $0.61$~bit/s/Hz.
    \end{alertblock}
    \begin{exampleblock}{Why this does not contradict the SISO result}
        A lens can increase received power at selected angles without
        improving channel orthogonality or the number of independent modes.
    \end{exampleblock}
\end{frame}

\subsection{Limitations and Next Steps}

\begin{frame}[t,shrink=22]{Main Limitations and Uncertainty}
    \begin{columns}[T,onlytextwidth]
        \begin{column}{0.48\textwidth}
            \begin{itemize}
                \item Only 20 drops
                \item 100,000 rather than 1,000,000 path samples
                \item Isotropic Monte Carlo TX pattern
                \item 10~dBm applied independently to every TX
                \item Ideal Shannon-bound capacity
            \end{itemize}
        \end{column}
        \begin{column}{0.48\textwidth}
            \begin{itemize}
                % \item Same patch CSV for all no-lens angle labels
                \item Lens-angle orientation and mapping require verification
                \item Virtual, not physical, $9\times9$ MIMO
                % \item LLVM JIT initialization warning in notebook output
            \end{itemize}
        \end{column}
    \end{columns}
    \begin{alertblock}{Fair-power comparison}
        For fixed total transmitted power, each TX should be reduced by
        $10\log_{10}(9)\approx9.54$~dB.
    \end{alertblock}
\end{frame}

% \begin{frame}[t,shrink=20]{Recommended Follow-Up Simulations}
%     \begin{itemize}
%         \item Increase to 100--500 drops and 1,000,000 path samples/source.
%         \item Export full-precision results to CSV or Parquet.
%         \item Report paired bootstrap confidence intervals for every angle.
%         \item Add outage thresholds at 5, 6, and 7~bit/s/Hz.
%         \item Repeat the comparison under a fixed total power budget.
%         \item Separate isotropic-TX and HFSS patch-TX experiments.
%         \item Verify boresight, angle signs, RX orientation, and HFSS mapping.
%         \item For physical MIMO, model simultaneously active RX elements with
%         an appropriate power-allocation and noise model.
%     \end{itemize}
% \end{frame}

\subsection{Conclusion}

\begin{frame}[t,shrink=15]{Monte Carlo Conclusions}
    \begin{columns}[T,onlytextwidth]
        \begin{column}{0.48\textwidth}
            \begin{exampleblock}{Promising evidence}
                \begin{itemize}
                    \item Median capacity: $+5.2\%$
                    \item Median throughput: $+5.1\%$
                    \item Median delay spread: $-4.8\%$
                    \item Best angle: $-30^\circ$
                \end{itemize}
            \end{exampleblock}
        \end{column}
        \begin{column}{0.48\textwidth}
            \begin{alertblock}{Important qualification}
                \begin{itemize}
                    \item Lower 5th-percentile capacity
                    \item $+60^\circ$ underperforms no lens
                    \item No virtual-MIMO conditioning improvement
                    \item Strong angle and pattern dependence
                \end{itemize}
            \end{alertblock}
        \end{column}
    \end{columns}
    \begin{block}{Engineering interpretation}
        The lens behaves as a direction-selective gain mechanism. It can
        improve typical performance when aligned, but it can be detrimental
        under angular mismatch. Final design decisions require more drops,
        fair total power, verified pattern mapping, and physical validation.
    \end{block}
\end{frame}
